\documentclass[11pt,a4paper]{article}
\usepackage[a4paper, margin=1.5cm]{geometry}
\usepackage{needspace}
\usepackage[nobottomtitles]{titlesec}
\usepackage{multicol}
\usepackage{xcolor}
\usepackage{wrapfig}
\usepackage{fp}
\usepackage{xfp}
\usepackage{pgfmath}
\usepackage{numprint}
\usepackage{siunitx}
\usepackage{enumitem}
\usepackage[version=4]{mhchem}
\usepackage{chemfig}
\usepackage{mol2chemfig}
\usepackage{WriteOnGrid}
\usepackage{frcursive}
\usepackage{csvsimple}%
\usepackage[francais,bloc]{automultiplechoice}

%\sisetup{locale = FR,detect-all,quotient-mode=fraction,
%input-product=*, list-final-separator = { \translate{et} },
%list-pair-separator = { \translate{et} },range-phrase = { \translate{à} },
%separate-uncertainty = true,group-minimum-digits=3}

%%% Force show explain during dev
%\newif\ifshowexplain
%\showexplaintrue  % Active l'affichage des explications
%\let\oldexplain\explain
%\renewcommand{\explain}[1]{%
%    \oldexplain{#1}% Conserve la définition originale pour AMC
%    \ifshowexplain
%        \par\noindent\textit{Explication : }#1\par
%    \fi
%}

\DeclareSIUnit{\nothing}{\relax}

\def\yseeda{12}
\def\yseedb{10}
\FPseed\yseeda
\pgfmathsetseed{\yseedb}

%\baremeDefautS{e=0,v=0,b=1,m=-.25}
\baremeDefautS{e=0,v=0,b=1,m=-.25}
\baremeDefautM{formula=(NBC>0 ? 1/NB : 0)}
%\baremeDefautM{formula=(NBC>0 ? (NMC==0 ? 1:0):0)}
%\AMCformula{questionmult}{scoring=v=1}

\graphicspath{ {./images/} }

\newenvironment{reponsesd}{
    \begin{multicols}{2}
    \begin{reponses} }{
    \end{reponses}
    \end{multicols}
}

% Sauvegarde l'environnement original
\let\oldEnvQuadrillage\EnvQuadrillage
\let\endoldEnvQuadrillage\endEnvQuadrillage

% Redéfinit l'environnement pour ajouter un espace après
\renewenvironment{EnvQuadrillage}[1][]{%
    \oldEnvQuadrillage[#1]%
}{%
    \endoldEnvQuadrillage%
    \vspace{0.5cm}% Ajoute un espace de 0,5 cm après chaque quadrillage
}

\renewenvironment{EnvQuadrillage}[1][]{}

% Sauvegarde l'ancienne définition de reponseshoriz
%\let\oldreponseshoriz\reponseshoriz
%\let\endoldreponseshoriz\endreponseshoriz

% Redéfinit reponseshoriz avec une taille de police réduite
%\renewenvironment{reponseshoriz}
%{
%    \oldreponseshoriz
%    \scriptsize % ou \small, \scriptsize, etc.
%}
%{
%    \endoldreponseshoriz
%}

\setlength{\columnseprule}{1pt}
\def\columnseprulecolor{\color{lightgray}}%

\let\oexplain\explain
\renewcommand{\explain}[1]{\oexplain{\textcolor{red}{#1}}}

\titleformat{\section}
  {\centering\hrule\vspace{2mm}}
  {\thesection}
  {1em}
  {}
  [\vspace{1mm}\hrule]

\titleformat{\subsection}
  {\em}
  {\thesubsection}
  {1em}
  {}

\newcommand{\sujet}{
\exemplaire{1}{%

\AMCsetFoot{\niveau{} \classe{} -- \prenom{}~\nom{} -- \thepage}
\AMCopenOpts{boxframerule=0pt,lineheight=0pt,width=0pt,boxframerulecol=white,
    dots=false,framerule=0pt,framerulecol=white}

%%% debut de l'en-tête des copies~:
\begin{center}
\vfill
\noindent{\large \bf Classe de \niveau{}°\classe{}}

\vspace*{3mm}

{\Large\bf Évaluation sur table STL - Physique-Chimie\\
04/03/2026}

\vfill
\namefield{\noindent{}\fbox{\vspace*{3mm}
         \Huge\bf\prenom{}~\nom{}\normalsize{}%
         \vspace*{3mm}
      }}
\end{center}
\vfill

\begin{center}
\textbf{Durée~: une heure.}
\vspace*{5mm}

  Aucun document n'est autorisé.
  L'usage de la calculatrice est interdit.

  Les questions faisant apparaître le symbole \multiSymbole{} peuvent
  présenter une ou plusieurs bonnes réponses. Les autres ont
  une unique bonne réponse.

  Des points négatifs seront affectés aux mauvaises réponses, \textbf{uniquement pour les QCM}. Pour les
  questions ouvertes (où vous écrivez la réponses dans les quadrillages), aucun point négatif ne sera appliqué.

  \vspace*{5mm}

  \textbf{IMPORTANT: Utilisez un crayon de papier bien noir pour cocher les cases, et une gomme
  pour effacer délicatement en cas d'erreur. Ne raturez pas les cases.
  Si vous effacez le pourtour de la case, ne le redessinez pas!}

  \vspace*{5mm}

  \textbf{IMPORTANT: Ne pas écrire dans les cases grisées!}

  \vspace*{5mm}

  Ne pas faire comme ceci (raturé, pas centré, trop pâle, ...):\\
  \includegraphics[width=4cm]{checkbox_bad.png}

  Mais comme cela (bien centré, bien foncé):\\
  \includegraphics[width=2cm]{checkbox_good.png}

  Si vous faites une erreur et que vous ne pouvez effacer, raturez la case bien ostensiblement.

\end{center}
\vspace{1ex}
\vfill
\pagebreak
%%% fin de l'en-tête

\restituegroupe{groupes}


\AMCassociation{\id}

	  } % End \exemplaire{1}{%
} % End \newcommand{\sujet}{


%%%%§§§§§§§§§§§§§§§§§§§§§§§§§§§§§§§§§
\newcommand{\allq}{}
\newcommand{\nball}[1]{#1}

\begin{document}
%%%Options
\AMCrandomseed{\yseeda}

%%\def\AMCformQuestion#1{\textbf{Question #1~:}}

\setdefaultgroupmode{withoutreplacement}
%%% Fin Options

%%% elements
\input{q_01_simplenum.tex}
\input{q_02_power10.tex}
\input{q_logs.tex}
\input{q_premdeg.tex}
\input{q_03_psi.tex}

% Composition d'un atome
\element{atoms}{
  \begin{question}{atoms1}
    Quelles sont les deux principales parties constitutives d'un atome~?
    \begin{reponses}
      \mauvaise{Électrons et photons}
      \mauvaise{Neutrons et protons}
      \mauvaise{Protons et électrons}
      \bonne{Noyau et électrons}
    \end{reponses}
    \explain{Un atome est constitué d'un noyau central et d'électrons en orbite autour de celui-ci.}
  \end{question}
}

\element{atoms}{
  \begin{question}{atoms2}
    Qu'est-ce qu'un nucléon~?
    \begin{reponses}
      \mauvaise{Un électron}
      \mauvaise{Un photon}
      \mauvaise{Une particule virtuelle}
      \bonne{Un proton ou un neutron}
    \end{reponses}
    \explain{Les nucléons sont les constituants du noyau atomique : protons et neutrons.}
  \end{question}
}

% Composition des noyaux
\element{atoms}{
  \begin{question}{atoms3}
    Quelles particules constituent principalement un noyau atomique~?
    \begin{reponses}
      \mauvaise{Électrons et neutrons}
      \mauvaise{Protons et photons}
      \mauvaise{Neutrons et photons}
      \bonne{Protons et neutrons}
    \end{reponses}
    \explain{Le noyau atomique est composé de protons et de neutrons.}
  \end{question}
}

\input{q_05_symchim.tex}

\input{q_alcans.tex}
\input{q_alcools.tex}
\input{q_groups.tex}
\input{q_nomencla.tex}
\input{q_isors.tex}
\input{q_isoze.tex}

\input{q_defacidebase.tex}
\input{q_phconcentration.tex}
\input{q_pkaforce.tex}
\input{q_tampons.tex}

\element{nox}{
    \begin{question}{nox1}
        Quel est le nombre d'oxydation de l'élément carbone dans \ce{CO2}~?
        \begin{reponses}
            \mauvaise{+1}
            \mauvaise{-2}
            \mauvaise{0}
            \bonne{+4}
        \end{reponses}
        \explain{Le nombre d'oxydation de l'oxygène est -II. Dans \ce{CO2}, la molécule est neutre, donc $ 2 \times (-II) + \text{no(C)} = 0 $. On en déduit que $\text{no(C)} = +IV$.}
    \end{question}
}

\element{nox}{
    \begin{question}{nox2}
        Quel est le nombre d'oxydation de l'élément azote dans \ce{NO3-}~?
        \begin{reponses}
            \mauvaise{+1}
            \mauvaise{+3}
            \mauvaise{-3}
            \bonne{+5}
        \end{reponses}
        \explain{Le nombre d'oxydation de l'oxygène est -II. L'ion \ce{NO3-} a une charge totale de -I, donc $ 3 \times (-II) + \text{no(N)} = -I $. On en déduit que $\text{no(N)} = +V$.}
    \end{question}
}

\element{nox}{
    \begin{question}{nox3}
        Quel est le nombre d'oxydation de l'élément soufre dans \ce{SO4^{2-}}~?
        \begin{reponses}
            \mauvaise{+2}
            \mauvaise{+4}
            \mauvaise{-2}
            \bonne{+6}
        \end{reponses}
        \explain{Le nombre d'oxydation de l'oxygène est -II. L'ion \ce{SO4^{2-}} a une charge totale de -II, donc $ 4 \times (-II) + \text{no(S)} = -II $. On en déduit que $\text{no(S)} = +VI$.}
    \end{question}
}

\element{nox}{
    \begin{question}{nox4}
        Quel est le nombre d'oxydation de l'élément manganèse dans \ce{MnO4-}~?
        \begin{reponses}
            \mauvaise{+3}
            \mauvaise{+5}
            \mauvaise{+1}
            \bonne{+7}
        \end{reponses}
        \explain{Le nombre d'oxydation de l'oxygène est -II. L'ion \ce{MnO4-} a une charge totale de -I, donc $ 4 \times (-II) + \text{no(Mn)} = -I $. On en déduit que $\text{no(Mn)} = +VII$.}
    \end{question}
}

\element{nox}{
    \begin{question}{nox5}
        Quel est le nombre d'oxydation de l'élément fer dans \ce{Fe2O3}~?
        \begin{reponses}
            \mauvaise{+1}
            \mauvaise{+4}
            \mauvaise{0}
            \bonne{+3}
        \end{reponses}
        \explain{Le nombre d'oxydation de l'oxygène est -II. Dans \ce{Fe2O3}, la molécule est neutre, donc $ 3 \times (-II) + 2 \times \text{no(Fe)} = 0 $. On en déduit que $\text{no(Fe)} = +III$.}
    \end{question}
}

\element{nox}{
    \begin{question}{nox6}
        Quel est le nombre d'oxydation de l'élément chlore dans \ce{ClO-}~?
        \begin{reponses}
            \mauvaise{-1}
            \mauvaise{+3}
            \mauvaise{0}
            \bonne{+1}
        \end{reponses}
        \explain{Le nombre d'oxydation de l'oxygène est -II. L'ion \ce{ClO-} a une charge totale de -I, donc $-II + \text{no(Cl)} = -I$. On en déduit que $\text{no(Cl)} = +I$.}
    \end{question}
}

\element{nox}{
    \begin{question}{nox7}
        Quel est le nombre d'oxydation de l'élément chrome dans \ce{Cr2O7^{2-}}~?
        \begin{reponses}
            \mauvaise{+2}
            \mauvaise{+4}
            \mauvaise{+5}
            \bonne{+6}
        \end{reponses}
        \explain{Le nombre d'oxydation de l'oxygène est -II. L'ion \ce{Cr2O7^{2-}} a une charge totale de -II, donc $ 7 \times (-II) + 2 \times \text{no(Cr)} = -II $. On en déduit que $\text{no(Cr)} = +VI$.}
    \end{question}
}

\element{nox}{
    \begin{question}{nox8}
        Quel est le nombre d'oxydation de l'élément hydrogène dans \ce{H2O}~?
        \begin{reponses}
            \mauvaise{0}
            \mauvaise{-2}
            \mauvaise{+2}
            \bonne{+1}
        \end{reponses}
        \explain{Le nombre d'oxydation de l'oxygène est -II. Dans \ce{H2O}, la molécule est neutre, donc $ 2 \times \text{no(H)} + (-II) = 0 $. On en déduit que $\text{no(H)} = +I$.}
    \end{question}
}


%%% fin des elements
\element{groupes}{
\section{Rappels simples sur des notions mathématiques et physiques}

\emph{Remarque: Le logarithme décimal est noté $\log_{10}$. Le logarithme naturel est noté $\ln$.}

\begin{multicols}{2}
\insertgroup[2]{simplenum}
\insertgroup[2]{power10}
\insertgroup[2]{logs}
\insertgroup[2]{premdeg}
\insertgroup[2]{psi}
\insertgroup[1]{atoms}
\insertgroup[1]{symchim}
\end{multicols}

\section{Nomenclature des alcanes}
\begin{multicols}{2}
\insertgroup[1]{falcans}
\insertgroup[1]{alcans}
\insertgroup[2]{alcools}
\insertgroup[1]{groups}
\insertgroup[1]{nomencla}
\insertgroup[1]{isors}
\insertgroup[1]{isoze}
\end{multicols}

\section{Couples acides/bases}
\begin{multicols}{2}
\insertgroup[1]{defacidebase}
\insertgroup[1]{phconcentration}
\insertgroup[1]{pkaforce}
\insertgroup[1]{tampons}
\end{multicols}

\section{Oxydo-réduction}

\emph{Rappel: Le nombre d'oxydation de l'hydrogène est $I$, celui de l'oxygène est $-II$.}

\begin{multicols}{2}
\insertgroup[4]{nox}
\end{multicols}

% Premier tirage (cinétique)
\pgfmathrandom{}
\let\yranOne\pgfmathresult
\pgfmathparse{\yranOne > 0.5 ? 1 : 0}
\let\fileChoiceOne\pgfmathresult

\ifnum\fileChoiceOne=1
    %\section{Saponification de l'éthanoate d'éthyle}
\textbf{Source du sujet :} Baccalauréat 2022.

\subsection{Introduction}
On souhaite déterminer le temps de demi-réaction de la saponification de l'éthanoate d'éthyle par deux méthodes : un suivi cinétique pH-métrique (partie A) puis l'exploitation mathématique de la loi de vitesse (partie B). L'équation de la réaction modélisant cette transformation chimique est :

\begin{figure}[h]
    \centering
% pubchem 8857
\chemfig{H_3C-[:30]C(=[:90]O)-[:330]O-[:30]CH_2-[:330]CH_3}
\caption{Éthanoate d'éthyle}
\end{figure}

\[
\ce{CH3COOCH2CH3 + HO^- -> CH3COO^- + CH3CH2OH}
\]

\subsection{Données}
\begin{itemize}
    \item Masse volumique de l'éthanoate d'éthyle : \(\rho_{\text{é.é.}} = 0,925~\mathrm{g/mL}\).
    \item Masse molaire de l'éthanoate d'éthyle : \(M_{\text{é.é.}} = 88,0~\mathrm{g/mol}\).
    \item Constante d'équilibre de l'autoprotolyse de l'eau : \(K_{\mathrm{e}} = 10^{-14}\).
\end{itemize}

\subsection{Partie A : Suivi pH-métrique de la saponification (2 points)}

Le suivi est réalisé en mettant en œuvre le protocole suivant :
\begin{itemize}
    \item Dans un bécher, verser :
    \begin{itemize}
        \item 10,0 mL de soude à 0,10 mol/L ;
        \item environ 35 mL d'eau ;
        \item environ 15 mL d'éthanol ;
        \item 2,0 mL d'éthanoate d'éthyle.
    \end{itemize}
    \item Déclencher l'acquisition informatique de la mesure du pH toutes les secondes. Le déclenchement définit l'instant initial \(t = 0~\mathrm{s}\).
\end{itemize}

\begin{questionmult}{saponif1}
    Déterminer les valeurs des quantités de matière initiales d'éthanoate d'éthyle et de soude.
    \begin{EnvQuadrillage}[NbCarreaux=10x4,Grille=Seyes,Marge=1]
    \end{EnvQuadrillage}
    \AMCOpen{}{
        \mauvaise[F]{NA}\scoring{b=0}
        \bonne[A1]{Calculs corrects}\scoring{b=1}
        \bonne[A2]{Résultats corrects}\scoring{b=1}
    }
    \explain{
    \begin{itemize}
        \item Pour la soude : \(n_{\text{HO}^-} = C \times V = 0,10~\mathrm{mol/L} \times 0,010~\mathrm{L} = 1,0 \times 10^{-3}~\mathrm{mol}\).
        \item Pour l'éthanoate d'éthyle : \(m = \rho \times V = 0,925~\mathrm{g/mL} \times 2,0~\mathrm{mL} = 1,85~\mathrm{g}\).
        \item Quantité de matière : \(n = \frac{m}{M} = \frac{1,85~\mathrm{g}}{88,0~\mathrm{g/mol}} = 2,10 \times 10^{-2}~\mathrm{mol}\).
    \end{itemize}
    }
\end{questionmult}

\begin{questionmult}{saponif2}
    Écrire l'équation de la réaction d'autoprotolyse de l'eau.
    \begin{EnvQuadrillage}[NbCarreaux=10x2,Grille=Seyes,Marge=1]
    \end{EnvQuadrillage}
    \AMCOpen{}{
        \mauvaise[F]{NA}\scoring{b=0}
        \bonne[A1]{Équation correcte}\scoring{b=1}
    }
    \explain{
    L'équation de l'autoprotolyse de l'eau est :
    \[
    \ce{2 H2O <=> H3O^+ + HO^-}
    \]
    }
\end{questionmult}

\begin{questionmult}{saponif3}
    Donner l'expression de la constante d'équilibre \(K_{\mathrm{e}}\) de la réaction d'autoprotolyse de l'eau en fonction des concentrations à l'équilibre.
    \begin{EnvQuadrillage}[NbCarreaux=10x2,Grille=Seyes,Marge=1]
    \end{EnvQuadrillage}
    \AMCOpen{}{
        \mauvaise[F]{NA}\scoring{b=0}
        \bonne[A1]{Expression correcte}\scoring{b=1}
    }
    \explain{
    La constante d'équilibre \(K_{\mathrm{e}}\) est donnée par :
    \[
    K_{\mathrm{e}} = [\ce{H3O^+}] \times [\ce{HO^-}]
    \]
    }
\end{questionmult}

\subsection{Partie B : Exploitation mathématique de la loi de vitesse (2 points)}

On considère que la concentration en éthanoate d'éthyle reste constante et égale à sa valeur initiale tout au long de la transformation chimique. La loi de vitesse de réaction peut alors s'écrire comme celle d'une loi d'ordre 1 : \(v = k_1 \times [\ce{HO^-}]\), où \(k_1\) est la constante de vitesse apparente.

\begin{questionmult}{saponif4}
    Montrer que la réaction est d'ordre 1 par rapport aux ions hydroxyde.
    \begin{EnvQuadrillage}[NbCarreaux=10x4,Grille=Seyes,Marge=1]
    \end{EnvQuadrillage}
    \AMCOpen{}{
        \mauvaise[F]{NA}\scoring{b=0}
        \bonne[A1]{Raisonnement correct}\scoring{b=1}
        \bonne[A2]{Conclusion correcte}\scoring{b=1}
    }
    \explain{
    La loi de vitesse \(v = k_1 \times [\ce{HO^-}]\) montre que la vitesse de réaction est directement proportionnelle à la concentration en ions hydroxyde. Cela indique que la réaction est d'ordre 1 par rapport aux ions hydroxyde.
    }
\end{questionmult}


\else
    %\section{Dépôt chimique en phase vapeur du carbure de silicium}
\textbf{Source du sujet~:} Baccalauréat STL Biotechnologies, Centres étrangers – 18 mai 2022.

\subsection{Introduction}
Le carbure de silicium, de formule SiC, a été découvert par Jons Jacob Berzelius en 1824 lors d'une expérience pour synthétiser du diamant. Il est devenu un matériau incontournable pour la fabrication d'instruments d'optique. Par exemple, il a été utilisé pour garantir la stabilité thermomécanique du télescope spatial infrarouge Herschel, développé par l'agence spatiale européenne et lancé en 2009. En particulier, la face optique des miroirs peut être revêtue de carbure de silicium par dépôt chimique en phase vapeur (ou CVD pour l'anglais « chemical vapor deposition ») afin de masquer toute porosité résiduelle et obtenir une surface polissable parfaite.

Dans ce procédé, un solide inerte servant de support est exposé à une ou plusieurs espèces chimiques en phase gazeuse qui se décomposent à sa surface pour former le matériau désiré. Parmi celles-ci, le méthyltrichlorosilane de formule $\ce{CH3SiCl3}$ est très souvent choisi. Par la suite, pour des raisons de simplification, il sera noté MTS.

On considère une enceinte vide, de volume constant, thermostatée à la température $T_{2} = 1200$ K, dans laquelle, au temps $t = 0$ min, on introduit une certaine quantité de MTS.

À cette température, la transformation permettant la formation de carbure de silicium peut être considérée comme totale.

L'équation de la réaction modélisant la transformation chimique au cours de laquelle le MTS se décompose est la suivante :
$$
\ce{CH3SiCl3(g) -> SiC(s) + 3HCl(g)}
$$

On suit par un procédé adapté l'évolution de la concentration en MTS au cours du temps. On obtient ainsi le graphe suivant :

\begin{figure}[h]
    \centering
    \includegraphics[width=0.8\linewidth]{ste7_silicium_graph.png}
    \caption{Concentration en MTS en fonction du temps}
    \label{fig:graphe_mts}
\end{figure}

\subsection{Questions}

\begin{questionmult}{sicQ0}
    Quel est l'ordre de la réaction de décomposition du MTS~? Justifier votre réponse.\\
    \begin{EnvQuadrillage}[NbCarreaux=18x4,Grille=Seyes,Marge=1]
    \end{EnvQuadrillage}
    \AMCOpen{}{
        \mauvaise[F]{NA}\scoring{b=0}
        \bonne[A1]{Ordre correct}\scoring{b=1}
        \bonne[A2]{Justification correcte}\scoring{b=1}
    }
    \explain{
    La réaction est d'ordre 1 par rapport au MTS, car la loi de vitesse est de la forme $v = k \times [\text{MTS}]$, et la concentration en MTS suit une loi exponentielle $C(t) = C_0 \cdot \mathrm{e}^{-kt}$.
    }
\end{questionmult}

\begin{questionmult}{sicQ1}
    Déterminer le temps de demi-réaction $t_{1/2}$ dans ces conditions expérimentales en expliquant votre démarche.\\
    \begin{EnvQuadrillage}[NbCarreaux=18x4,Grille=Seyes,Marge=1]
    \end{EnvQuadrillage}
    \AMCOpen{}{
        \mauvaise[F]{NA}\scoring{b=0}
        \bonne[A1]{Méthode graphique correcte}\scoring{b=1}
        \bonne[A2]{Valeur correcte}\scoring{b=1}
    }
    \explain{
    Le temps de demi-réaction $t_{1/2}$ est déterminé graphiquement en repérant le point où la concentration en MTS est égale à la moitié de la concentration initiale.
    }
\end{questionmult}

\begin{questionmult}{sicQ2}
    On rappelle que $t_{1/2} = \frac{\ln 2}{k}$ avec $k$ la constante de vitesse de la réaction. Déterminer la valeur de $k$ dont on précisera l'unité.\\
    \begin{EnvQuadrillage}[NbCarreaux=10x2,Grille=Seyes,Marge=1]
    \end{EnvQuadrillage}
    \AMCOpen{}{
        \mauvaise[F]{NA}\scoring{b=0}
        \bonne[A1]{Expression correcte}\scoring{b=1}
        \bonne[A2]{Valeur correcte}\scoring{b=1}
    }
    \explain{
    La constante de vitesse $k$ est calculée à partir de la relation $t_{1/2} = \frac{\ln 2}{k}$. L'unité de $k$ est $\min^{-1}$.
    }
\end{questionmult}

\begin{questionmult}{sicQ3}
    Déterminer graphiquement la valeur de la vitesse de disparition du MTS à l'instant $t = 10$ min.\\
    \begin{EnvQuadrillage}[NbCarreaux=10x2,Grille=Seyes,Marge=1]
    \end{EnvQuadrillage}
    \AMCOpen{}{
        \mauvaise[F]{NA}\scoring{b=0}
        \bonne[A1]{Méthode correcte}\scoring{b=1}
        \bonne[A2]{Valeur correcte}\scoring{b=1}
    }
    \explain{
    La vitesse de disparition du MTS à $t = 10$ min est déterminée par la pente de la tangente à la courbe $C(t)$ en $t = 10$ min.
    }
\end{questionmult}

\begin{questionmult}{sicQ4}
    La vitesse de disparition du MTS est de $1,0 \times 10^{-2}$ mol.L$^{-1}$.min$^{-1}$ à $t = 1$ min et de $3,5 \times 10^{-3}$ mol.L$^{-1}$.min$^{-1}$ à $t = 30$ min. Conclure en discutant de l'évolution au cours du temps de la vitesse de disparition du MTS lorsque la concentration évolue.\\
    \begin{EnvQuadrillage}[NbCarreaux=18x4,Grille=Seyes,Marge=1]
    \end{EnvQuadrillage}
    \AMCOpen{}{
        \mauvaise[F]{NA}\scoring{b=0}
        \bonne[A1]{Comparaison correcte}\scoring{b=1}
        \bonne[A2]{Conclusion correcte}\scoring{b=1}
    }
    \explain{
    Les vitesses de disparition du MTS sont de $1,0 \times 10^{-2}$ mol.L$^{-1}$.min$^{-1}$ à $t = 1$ min et de $3,5 \times 10^{-3}$ mol.L$^{-1}$.min$^{-1}$ à $t = 30$ min. La vitesse de disparition diminue au cours du temps.
    }
\end{questionmult}

\subsection{Modélisation mathématique}

On modélise la concentration en MTS exprimée en mol.L$^{-1}$ en fonction du temps $t$ exprimé en minute, par la fonction $C$, définie sur l'intervalle [0; 50] par :
$$
C(t) = 0,30 \cdot \mathrm{e}^{-0,035t}
$$

\begin{questionmult}{sicQ5}
    On note $C'$ la fonction dérivée de la fonction $C$ sur l'intervalle [0; 50]. Déterminer l'expression de $C'(t)$ pour $t$ appartenant à [0; 50].
    \begin{EnvQuadrillage}[NbCarreaux=10x4,Grille=Seyes,Marge=1]
    \end{EnvQuadrillage}
    \AMCOpen{}{
        \mauvaise[F]{NA}\scoring{b=0}
        \bonne[A1]{Dérivation correcte}\scoring{b=1}
    }
    \explain{
    La dérivée de $C(t) = 0,30 \cdot \mathrm{e}^{-0,035t}$ est $C'(t) = -0,0105 \mathrm{e}^{-0,035t}$.
    }
\end{questionmult}

\begin{questionmult}{sicQ6}
    On rappelle que la vitesse de disparition de MTS est égale à l'opposé de la fonction dérivée $C'$. On note $C''$ la fonction dérivée de $C'$. On admet que $C''(t) = 3,675 \cdot 10^{-4} \mathrm{e}^{-0,035t}$ pour $t$ appartenant à [0; 50]. Étudier le sens de variation de la vitesse de réaction au cours du temps.\\
    \begin{EnvQuadrillage}[NbCarreaux=18x6,Grille=Seyes,Marge=1]
    \end{EnvQuadrillage}
    \AMCOpen{}{
        \mauvaise[F]{NA}\scoring{b=0}
        \bonne[A1]{Analyse correcte}\scoring{b=1}
        \bonne[A2]{Conclusion correcte}\scoring{b=1}
    }
    \explain{
    La dérivée seconde $C''(t) = 3,675 \times 10^{-4} \mathrm{e}^{-0,035t}$ est positive sur [0; 50]. Donc $C'(t)$ est croissante, et la vitesse de disparition (opposée de $C'(t)$) est décroissante.
    }
\end{questionmult}

\begin{questionmult}{sicQ7}
    On considère que la transformation chimique de décomposition de MTS peut être stoppée lorsqu'il ne reste que $10\%$ de la concentration initiale de MTS. Déterminer l'instant $t$ à partir duquel la transformation chimique peut être stoppée. On donnera la valeur exacte, puis la valeur arrondie à la minute près.\\
    \begin{EnvQuadrillage}[NbCarreaux=14x6,Grille=Seyes,Marge=1]
    \end{EnvQuadrillage}
    \AMCOpen{}{
        \mauvaise[F]{NA}\scoring{b=0}
        \bonne[A1]{Résolution correcte}\scoring{b=1}
        \bonne[A2]{Valeur correcte}\scoring{b=1}
    }
    \explain{
    On résout $C(t) = 0,03$ :
    $$
    0,30 \cdot \mathrm{e}^{-0,035t} = 0,03 \implies t = \frac{\ln 10}{0,035} \approx 66 \text{ min}.
    $$
    }
\end{questionmult}

\fi

\section{Saponification de l'éthanoate d'éthyle}
\textbf{Source du sujet :} Baccalauréat 2022.

\subsection{Introduction}
On souhaite déterminer le temps de demi-réaction de la saponification de l'éthanoate d'éthyle par deux méthodes : un suivi cinétique pH-métrique (partie A) puis l'exploitation mathématique de la loi de vitesse (partie B). L'équation de la réaction modélisant cette transformation chimique est :

\begin{figure}[h]
    \centering
% pubchem 8857
\chemfig{H_3C-[:30]C(=[:90]O)-[:330]O-[:30]CH_2-[:330]CH_3}
\caption{Éthanoate d'éthyle}
\end{figure}

\[
\ce{CH3COOCH2CH3 + HO^- -> CH3COO^- + CH3CH2OH}
\]

\subsection{Données}
\begin{itemize}
    \item Masse volumique de l'éthanoate d'éthyle : \(\rho_{\text{é.é.}} = 0,925~\mathrm{g/mL}\).
    \item Masse molaire de l'éthanoate d'éthyle : \(M_{\text{é.é.}} = 88,0~\mathrm{g/mol}\).
    \item Constante d'équilibre de l'autoprotolyse de l'eau : \(K_{\mathrm{e}} = 10^{-14}\).
\end{itemize}

\subsection{Partie A : Suivi pH-métrique de la saponification (2 points)}

Le suivi est réalisé en mettant en œuvre le protocole suivant :
\begin{itemize}
    \item Dans un bécher, verser :
    \begin{itemize}
        \item 10,0 mL de soude à 0,10 mol/L ;
        \item environ 35 mL d'eau ;
        \item environ 15 mL d'éthanol ;
        \item 2,0 mL d'éthanoate d'éthyle.
    \end{itemize}
    \item Déclencher l'acquisition informatique de la mesure du pH toutes les secondes. Le déclenchement définit l'instant initial \(t = 0~\mathrm{s}\).
\end{itemize}

\begin{questionmult}{saponif1}
    Déterminer les valeurs des quantités de matière initiales d'éthanoate d'éthyle et de soude.
    \begin{EnvQuadrillage}[NbCarreaux=10x4,Grille=Seyes,Marge=1]
    \end{EnvQuadrillage}
    \AMCOpen{}{
        \mauvaise[F]{NA}\scoring{b=0}
        \bonne[A1]{Calculs corrects}\scoring{b=1}
        \bonne[A2]{Résultats corrects}\scoring{b=1}
    }
    \explain{
    \begin{itemize}
        \item Pour la soude : \(n_{\text{HO}^-} = C \times V = 0,10~\mathrm{mol/L} \times 0,010~\mathrm{L} = 1,0 \times 10^{-3}~\mathrm{mol}\).
        \item Pour l'éthanoate d'éthyle : \(m = \rho \times V = 0,925~\mathrm{g/mL} \times 2,0~\mathrm{mL} = 1,85~\mathrm{g}\).
        \item Quantité de matière : \(n = \frac{m}{M} = \frac{1,85~\mathrm{g}}{88,0~\mathrm{g/mol}} = 2,10 \times 10^{-2}~\mathrm{mol}\).
    \end{itemize}
    }
\end{questionmult}

\begin{questionmult}{saponif2}
    Écrire l'équation de la réaction d'autoprotolyse de l'eau.
    \begin{EnvQuadrillage}[NbCarreaux=10x2,Grille=Seyes,Marge=1]
    \end{EnvQuadrillage}
    \AMCOpen{}{
        \mauvaise[F]{NA}\scoring{b=0}
        \bonne[A1]{Équation correcte}\scoring{b=1}
    }
    \explain{
    L'équation de l'autoprotolyse de l'eau est :
    \[
    \ce{2 H2O <=> H3O^+ + HO^-}
    \]
    }
\end{questionmult}

\begin{questionmult}{saponif3}
    Donner l'expression de la constante d'équilibre \(K_{\mathrm{e}}\) de la réaction d'autoprotolyse de l'eau en fonction des concentrations à l'équilibre.
    \begin{EnvQuadrillage}[NbCarreaux=10x2,Grille=Seyes,Marge=1]
    \end{EnvQuadrillage}
    \AMCOpen{}{
        \mauvaise[F]{NA}\scoring{b=0}
        \bonne[A1]{Expression correcte}\scoring{b=1}
    }
    \explain{
    La constante d'équilibre \(K_{\mathrm{e}}\) est donnée par :
    \[
    K_{\mathrm{e}} = [\ce{H3O^+}] \times [\ce{HO^-}]
    \]
    }
\end{questionmult}

\subsection{Partie B : Exploitation mathématique de la loi de vitesse (2 points)}

On considère que la concentration en éthanoate d'éthyle reste constante et égale à sa valeur initiale tout au long de la transformation chimique. La loi de vitesse de réaction peut alors s'écrire comme celle d'une loi d'ordre 1 : \(v = k_1 \times [\ce{HO^-}]\), où \(k_1\) est la constante de vitesse apparente.

\begin{questionmult}{saponif4}
    Montrer que la réaction est d'ordre 1 par rapport aux ions hydroxyde.
    \begin{EnvQuadrillage}[NbCarreaux=10x4,Grille=Seyes,Marge=1]
    \end{EnvQuadrillage}
    \AMCOpen{}{
        \mauvaise[F]{NA}\scoring{b=0}
        \bonne[A1]{Raisonnement correct}\scoring{b=1}
        \bonne[A2]{Conclusion correcte}\scoring{b=1}
    }
    \explain{
    La loi de vitesse \(v = k_1 \times [\ce{HO^-}]\) montre que la vitesse de réaction est directement proportionnelle à la concentration en ions hydroxyde. Cela indique que la réaction est d'ordre 1 par rapport aux ions hydroxyde.
    }
\end{questionmult}


\section{Dépôt chimique en phase vapeur du carbure de silicium}
\textbf{Source du sujet~:} Baccalauréat STL Biotechnologies, Centres étrangers – 18 mai 2022.

\subsection{Introduction}
Le carbure de silicium, de formule SiC, a été découvert par Jons Jacob Berzelius en 1824 lors d'une expérience pour synthétiser du diamant. Il est devenu un matériau incontournable pour la fabrication d'instruments d'optique. Par exemple, il a été utilisé pour garantir la stabilité thermomécanique du télescope spatial infrarouge Herschel, développé par l'agence spatiale européenne et lancé en 2009. En particulier, la face optique des miroirs peut être revêtue de carbure de silicium par dépôt chimique en phase vapeur (ou CVD pour l'anglais « chemical vapor deposition ») afin de masquer toute porosité résiduelle et obtenir une surface polissable parfaite.

Dans ce procédé, un solide inerte servant de support est exposé à une ou plusieurs espèces chimiques en phase gazeuse qui se décomposent à sa surface pour former le matériau désiré. Parmi celles-ci, le méthyltrichlorosilane de formule $\ce{CH3SiCl3}$ est très souvent choisi. Par la suite, pour des raisons de simplification, il sera noté MTS.

On considère une enceinte vide, de volume constant, thermostatée à la température $T_{2} = 1200$ K, dans laquelle, au temps $t = 0$ min, on introduit une certaine quantité de MTS.

À cette température, la transformation permettant la formation de carbure de silicium peut être considérée comme totale.

L'équation de la réaction modélisant la transformation chimique au cours de laquelle le MTS se décompose est la suivante :
$$
\ce{CH3SiCl3(g) -> SiC(s) + 3HCl(g)}
$$

On suit par un procédé adapté l'évolution de la concentration en MTS au cours du temps. On obtient ainsi le graphe suivant :

\begin{figure}[h]
    \centering
    \includegraphics[width=0.8\linewidth]{ste7_silicium_graph.png}
    \caption{Concentration en MTS en fonction du temps}
    \label{fig:graphe_mts}
\end{figure}

\subsection{Questions}

\begin{questionmult}{sicQ0}
    Quel est l'ordre de la réaction de décomposition du MTS~? Justifier votre réponse.\\
    \begin{EnvQuadrillage}[NbCarreaux=18x4,Grille=Seyes,Marge=1]
    \end{EnvQuadrillage}
    \AMCOpen{}{
        \mauvaise[F]{NA}\scoring{b=0}
        \bonne[A1]{Ordre correct}\scoring{b=1}
        \bonne[A2]{Justification correcte}\scoring{b=1}
    }
    \explain{
    La réaction est d'ordre 1 par rapport au MTS, car la loi de vitesse est de la forme $v = k \times [\text{MTS}]$, et la concentration en MTS suit une loi exponentielle $C(t) = C_0 \cdot \mathrm{e}^{-kt}$.
    }
\end{questionmult}

\begin{questionmult}{sicQ1}
    Déterminer le temps de demi-réaction $t_{1/2}$ dans ces conditions expérimentales en expliquant votre démarche.\\
    \begin{EnvQuadrillage}[NbCarreaux=18x4,Grille=Seyes,Marge=1]
    \end{EnvQuadrillage}
    \AMCOpen{}{
        \mauvaise[F]{NA}\scoring{b=0}
        \bonne[A1]{Méthode graphique correcte}\scoring{b=1}
        \bonne[A2]{Valeur correcte}\scoring{b=1}
    }
    \explain{
    Le temps de demi-réaction $t_{1/2}$ est déterminé graphiquement en repérant le point où la concentration en MTS est égale à la moitié de la concentration initiale.
    }
\end{questionmult}

\begin{questionmult}{sicQ2}
    On rappelle que $t_{1/2} = \frac{\ln 2}{k}$ avec $k$ la constante de vitesse de la réaction. Déterminer la valeur de $k$ dont on précisera l'unité.\\
    \begin{EnvQuadrillage}[NbCarreaux=10x2,Grille=Seyes,Marge=1]
    \end{EnvQuadrillage}
    \AMCOpen{}{
        \mauvaise[F]{NA}\scoring{b=0}
        \bonne[A1]{Expression correcte}\scoring{b=1}
        \bonne[A2]{Valeur correcte}\scoring{b=1}
    }
    \explain{
    La constante de vitesse $k$ est calculée à partir de la relation $t_{1/2} = \frac{\ln 2}{k}$. L'unité de $k$ est $\min^{-1}$.
    }
\end{questionmult}

\begin{questionmult}{sicQ3}
    Déterminer graphiquement la valeur de la vitesse de disparition du MTS à l'instant $t = 10$ min.\\
    \begin{EnvQuadrillage}[NbCarreaux=10x2,Grille=Seyes,Marge=1]
    \end{EnvQuadrillage}
    \AMCOpen{}{
        \mauvaise[F]{NA}\scoring{b=0}
        \bonne[A1]{Méthode correcte}\scoring{b=1}
        \bonne[A2]{Valeur correcte}\scoring{b=1}
    }
    \explain{
    La vitesse de disparition du MTS à $t = 10$ min est déterminée par la pente de la tangente à la courbe $C(t)$ en $t = 10$ min.
    }
\end{questionmult}

\begin{questionmult}{sicQ4}
    La vitesse de disparition du MTS est de $1,0 \times 10^{-2}$ mol.L$^{-1}$.min$^{-1}$ à $t = 1$ min et de $3,5 \times 10^{-3}$ mol.L$^{-1}$.min$^{-1}$ à $t = 30$ min. Conclure en discutant de l'évolution au cours du temps de la vitesse de disparition du MTS lorsque la concentration évolue.\\
    \begin{EnvQuadrillage}[NbCarreaux=18x4,Grille=Seyes,Marge=1]
    \end{EnvQuadrillage}
    \AMCOpen{}{
        \mauvaise[F]{NA}\scoring{b=0}
        \bonne[A1]{Comparaison correcte}\scoring{b=1}
        \bonne[A2]{Conclusion correcte}\scoring{b=1}
    }
    \explain{
    Les vitesses de disparition du MTS sont de $1,0 \times 10^{-2}$ mol.L$^{-1}$.min$^{-1}$ à $t = 1$ min et de $3,5 \times 10^{-3}$ mol.L$^{-1}$.min$^{-1}$ à $t = 30$ min. La vitesse de disparition diminue au cours du temps.
    }
\end{questionmult}

\subsection{Modélisation mathématique}

On modélise la concentration en MTS exprimée en mol.L$^{-1}$ en fonction du temps $t$ exprimé en minute, par la fonction $C$, définie sur l'intervalle [0; 50] par :
$$
C(t) = 0,30 \cdot \mathrm{e}^{-0,035t}
$$

\begin{questionmult}{sicQ5}
    On note $C'$ la fonction dérivée de la fonction $C$ sur l'intervalle [0; 50]. Déterminer l'expression de $C'(t)$ pour $t$ appartenant à [0; 50].
    \begin{EnvQuadrillage}[NbCarreaux=10x4,Grille=Seyes,Marge=1]
    \end{EnvQuadrillage}
    \AMCOpen{}{
        \mauvaise[F]{NA}\scoring{b=0}
        \bonne[A1]{Dérivation correcte}\scoring{b=1}
    }
    \explain{
    La dérivée de $C(t) = 0,30 \cdot \mathrm{e}^{-0,035t}$ est $C'(t) = -0,0105 \mathrm{e}^{-0,035t}$.
    }
\end{questionmult}

\begin{questionmult}{sicQ6}
    On rappelle que la vitesse de disparition de MTS est égale à l'opposé de la fonction dérivée $C'$. On note $C''$ la fonction dérivée de $C'$. On admet que $C''(t) = 3,675 \cdot 10^{-4} \mathrm{e}^{-0,035t}$ pour $t$ appartenant à [0; 50]. Étudier le sens de variation de la vitesse de réaction au cours du temps.\\
    \begin{EnvQuadrillage}[NbCarreaux=18x6,Grille=Seyes,Marge=1]
    \end{EnvQuadrillage}
    \AMCOpen{}{
        \mauvaise[F]{NA}\scoring{b=0}
        \bonne[A1]{Analyse correcte}\scoring{b=1}
        \bonne[A2]{Conclusion correcte}\scoring{b=1}
    }
    \explain{
    La dérivée seconde $C''(t) = 3,675 \times 10^{-4} \mathrm{e}^{-0,035t}$ est positive sur [0; 50]. Donc $C'(t)$ est croissante, et la vitesse de disparition (opposée de $C'(t)$) est décroissante.
    }
\end{questionmult}

\begin{questionmult}{sicQ7}
    On considère que la transformation chimique de décomposition de MTS peut être stoppée lorsqu'il ne reste que $10\%$ de la concentration initiale de MTS. Déterminer l'instant $t$ à partir duquel la transformation chimique peut être stoppée. On donnera la valeur exacte, puis la valeur arrondie à la minute près.\\
    \begin{EnvQuadrillage}[NbCarreaux=14x6,Grille=Seyes,Marge=1]
    \end{EnvQuadrillage}
    \AMCOpen{}{
        \mauvaise[F]{NA}\scoring{b=0}
        \bonne[A1]{Résolution correcte}\scoring{b=1}
        \bonne[A2]{Valeur correcte}\scoring{b=1}
    }
    \explain{
    On résout $C(t) = 0,03$ :
    $$
    0,30 \cdot \mathrm{e}^{-0,035t} = 0,03 \implies t = \frac{\ln 10}{0,035} \approx 66 \text{ min}.
    $$
    }
\end{questionmult}


% Deuxième tirage (redox)
\pgfmathrandom{}
\let\yranTwo\pgfmathresult
\pgfmathparse{\yranTwo > 0.5 ? 1 : 0}
\let\fileChoiceTwo\pgfmathresult

\ifnum\fileChoiceTwo=1
    %\section{Comparaison des piles pour appareil auditif}
\textbf{Source du sujet~:} Baccalauréat STL Physique-Chimie, Centres étrangers – 2022.

\subsection{Introduction}
Pour alimenter un appareil auditif, deux types de piles boutons sont disponibles :
\begin{itemize}
    \item une pile bouton zinc-manganèse ;
    \item une pile bouton zinc-air.
\end{itemize}

Leurs caractéristiques sont les suivantes :

\begin{center}
\begin{tabular}{|c|c|c|c|c|}
\hline
Type de pile & Quantité d'électricité disponible & Masse de la pile & Tension & Intensité du courant \\
\hline
zinc-manganèse & à déterminer & 0,5 g & 1,5 V & 2,0 mA \\
\hline
zinc-air & 182 mA·h & 0,5 g & 1,5 V & 2,0 mA \\
\hline
\end{tabular}
\end{center}

L'objectif de l'exercice est de déterminer la pile à privilégier pour une plus grande autonomie de fonctionnement de l'appareil auditif.

\subsection{Données}
\begin{itemize}
    \item Masses molaires atomiques : $M(\mathrm{Zn}) = \SI{65}{\gram\per\mole}$ ; $M(\mathrm{Mn}) = \SI{55}{\gram\per\mole}$ ; $M(\mathrm{O}) = \SI{16}{\gram\per\mole}$ ; $M(\mathrm{H}) = \SI{1,0}{\gram\per\mole}$ ;
    \item Constante de Faraday : $F = \SI{9,65e4}{\coulomb\per\mole}$.
\end{itemize}

\subsection{Partie A~: Étude de la pile zinc-manganèse}
Dans une pile bouton zinc-manganèse, une électrode contient du zinc $\ce{Zn(s)}$ et l'autre du dioxyde de manganèse $\ce{MnO2(s)}$ et du carbone $\ce{C(s)}$. Le carbone n'intervient pas dans le couple de la demi-pile mais permet notamment la conduction du courant.\\
Afin de déterminer la polarité de la pile, on branche la borne COM d’un voltmètre à
l’électrode de carbone et la borne V à l’électrode de zinc, on mesure alors une tension
de \SI{-1.5}{\volt}.\\

\begin{questionmult}{pileA1}
    Déduire de la mesure de la tension à vide la polarité de chaque électrode en précisant le raisonnement.\\
    \begin{EnvQuadrillage}[NbCarreaux=18x3,Grille=Seyes,Marge=1]
    \end{EnvQuadrillage}
    \AMCOpen{}{
        \mauvaise[F]{NA}\scoring{b=0}
        \bonne[A1]{idOK}\scoring{b=1}
        \bonne[A2]{valOK}\scoring{b=1}
    }
    \explain{
    La tension mesurée est $-1,50\ \mathrm{V}$ lorsque la borne COM est branchée sur l'électrode de carbone et la borne V sur l'électrode de zinc. Cela indique que l'électrode de carbone est la borne positive et l'électrode de zinc la borne négative.
    }
\end{questionmult}

\begin{questionmult}{pileA2}
    Indiquer la polarité des bornes de la pile, le sens de déplacement des électrons ainsi que le sens du courant sur la figure.

    \includegraphics[width=0.6\linewidth]{ste7_pile_auditif_pile.png}

    \AMCOpen{}{
        \mauvaise[F]{NA}\scoring{b=0}
        \bonne[A1]{PolOK}\scoring{b=1}
        \bonne[A2]{SensOK}\scoring{b=1}
    }
    \explain{
    Les électrons se déplacent de l'électrode de zinc vers l'électrode de carbone, et le courant circule dans le sens opposé.
    }
\end{questionmult}

 \vspace{0.5cm}

Les demi-équations électrochimiques se produisant aux électrodes de la pile sont les suivantes :

$$
\ce{Zn(s) + 2 HO^{-}(aq) -> ZnO(s) + H2O(l) + 2 e^{-}}
$$

$$
\ce{MnO2(s) + H2O(l) + 1 e^{-} -> MnO2H(s) + HO^{-}(aq)}
$$

\begin{questionmult}{pileA3}
    Associer chaque demi-équation électrochimique à l'électrode correspondante et préciser si elle est le siège d'une oxydation ou d'une réduction.\\
    \begin{EnvQuadrillage}[NbCarreaux=18x4,Grille=Seyes,Marge=1]
    \end{EnvQuadrillage}
    \AMCOpen{}{
        \mauvaise{NA}\scoring{b=0}
        \bonne{assOK}\scoring{b=1}
        \bonne{redoxOK}\scoring{b=1}
    }
    \explain{
    La première demi-équation correspond à l'oxydation du zinc à l'anode. La deuxième demi-équation correspond à la réduction du dioxyde de manganèse à la cathode.
    }
\end{questionmult}

\begin{questionmult}{pileA4}
    Établir l'équation de la réaction modélisant le fonctionnement de la pile, en précisant la démarche.\\
    \begin{EnvQuadrillage}[NbCarreaux=18x8,Grille=Seyes,Marge=1]
    \end{EnvQuadrillage}
    \AMCOpen{}{
        \mauvaise{NA}\scoring{b=0}
        \bonne{idK}\scoring{b=1}
        \bonne{eqOK}\scoring{b=1}
    }
    \explain{
    L'équation de la réaction est obtenue en combinant les deux demi-équations, en tenant compte du nombre d'électrons échangés :
    $$
    \ce{2 Zn(s) + MnO2(s) + H2O(l) -> 2 ZnO(s) + MnO2H(s)}
    $$
    }
\end{questionmult}

\begin{questionmult}{pileA5}
    Sachant que la masse initiale de l'électrode en zinc est $m_{\mathrm{Zn}} = 55\ \mathrm{mg}$, déterminer la quantité de matière initiale de zinc $n_{\mathrm{Zn}}$ présente dans la pile.\\
    \begin{EnvQuadrillage}[NbCarreaux=18x3,Grille=Seyes,Marge=1]
    \end{EnvQuadrillage}
    \AMCOpen{}{
        \mauvaise[F]{NA}\scoring{b=0}
        \bonne[A1]{idOK}\scoring{b=1}
        \bonne[A2]{valOK}\scoring{b=1}
    }
    \explain{
    La quantité de matière initiale de zinc est donnée par $n_{\mathrm{Zn}} = \frac{m_{\mathrm{Zn}}}{M(\mathrm{Zn})} = \frac{55 \times 10^{-3}}{65} = 8,46 \times 10^{-4}\ \mathrm{mol}$.
    }
\end{questionmult}

\begin{questionmult}{pileA6}
    En admettant que le zinc est entièrement consommé et que le dioxyde de manganèse $\ce{MnO2(s)}$ est en excès, déduire la quantité de matière $n_{\mathrm{e}^{-}}$ d'électrons échangés au cours de la transformation chimique.\\
    \begin{EnvQuadrillage}[NbCarreaux=18x3,Grille=Seyes,Marge=1]
    \end{EnvQuadrillage}
    \AMCOpen{}{
        \mauvaise[F]{NA}\scoring{b=0}
        \bonne[A1]{idOK}\scoring{b=1}
        \bonne[A2]{valOK}\scoring{b=1}
    }
    \explain{
    Chaque mole de zinc libère 2 moles d'électrons. Donc, $n_{\mathrm{e}^{-}} = 2 \times n_{\mathrm{Zn}} = 1,69 \times 10^{-3}\ \mathrm{mol}$.
    }
\end{questionmult}

\begin{questionmult}{pileA7}
    Vérifier que la valeur de la quantité d'électricité disponible dans la pile est voisine de $Q_{1} = 164\ \mathrm{C}$.\\
    \begin{EnvQuadrillage}[NbCarreaux=10x2,Grille=Seyes,Marge=1]
    \end{EnvQuadrillage}
    \AMCOpen{}{
        \mauvaise[F]{NA}\scoring{b=0}
        \bonne[A1]{idOK}\scoring{b=1}
        \bonne[A2]{valOK}\scoring{b=1}
    }
    \explain{
    La quantité d'électricité est donnée par $Q = n_{\mathrm{e}^{-}} \times F = 1,69 \times 10^{-3} \times 9,65 \times 10^{4} = 163\ \mathrm{C}$.
    }
\end{questionmult}

\begin{questionmult}{pileA8}
    Calculer la durée de fonctionnement de la pile zinc-manganèse.\\
    \begin{EnvQuadrillage}[NbCarreaux=10x2,Grille=Seyes,Marge=1]
    \end{EnvQuadrillage}
    \AMCOpen{}{
        \mauvaise[F]{NA}\scoring{b=0}
        \bonne[A1]{idOK}\scoring{b=1}
        \bonne[A2]{valOK}\scoring{b=1}
    }
    \explain{
    La durée de fonctionnement est donnée par $t = \frac{Q}{I} = \frac{163}{2,0 \times 10^{-3}} = 81\,500\ \mathrm{s} = 22,6\ \mathrm{h}$.
    }
\end{questionmult}

\subsection{Partie B~: Étude de la pile zinc-air}
L'équation de la réaction modélisant le fonctionnement de la pile correspond à l'oxydation du zinc en oxyde de zinc par le dioxygène de l'air :

$$
\ce{2 Zn(s) + O2(g) -> 2 ZnO(s)}
$$

Un des deux couples redox mis en jeu est le couple $\ce{O2(g) / H2O(l)}$.

\begin{questionmult}{pileB1}
    Écrire l'équation de demi-réaction électrochimique correspondant à ce couple et justifier le caractère oxydant du dioxygène.\\
    \begin{EnvQuadrillage}[NbCarreaux=18x4,Grille=Seyes,Marge=1]
    \end{EnvQuadrillage}
    \AMCOpen{}{
        \mauvaise[F]{NA}\scoring{b=0}
        \bonne[A1]{Équation correcte}\scoring{b=1}
        \bonne[A2]{Justification correcte}\scoring{b=1}
    }
    \explain{
    La demi-équation est $\ce{O2(g) + 2 H2O(l) + 4 e^{-} -> 4 HO^{-}(aq)}$. Le dioxygène est un oxydant car il gagne des électrons.
    }
\end{questionmult}

\begin{questionmult}{pileB2}
    Calculer la durée de fonctionnement de la pile zinc-air en vous aidant de ses caractéristiques.\\
    \begin{EnvQuadrillage}[NbCarreaux=18x4,Grille=Seyes,Marge=1]
    \end{EnvQuadrillage}
    \AMCOpen{}{
        \mauvaise[F]{NA}\scoring{b=0}
        \bonne[A1]{idOK}\scoring{b=1}
        \bonne[A2]{valOK}\scoring{b=1}
    }
    \explain{
    La quantité d'électricité disponible est $Q = 182\ \mathrm{mA\cdot h} = 655,2\ \mathrm{C}$. La durée de fonctionnement est $t = \frac{Q}{I} = \frac{655,2}{2,0 \times 10^{-3}} = 327\,600\ \mathrm{s} = 91\ \mathrm{h}$.
    }
\end{questionmult}

\subsection{Partie C~: Choix de la pile à privilégier}
\begin{questionmult}{pileC1}
    À l'aide des parties A et B et des informations disponibles, indiquer le type de pile qu'il faudrait privilégier pour l'alimentation d'un appareil auditif. Justifier la réponse.\\
    \begin{EnvQuadrillage}[NbCarreaux=18x4,Grille=Seyes,Marge=1]
    \end{EnvQuadrillage}
    \AMCOpen{}{
        \mauvaise[F]{NA}\scoring{b=0}
        \bonne[A1]{id}\scoring{b=1}
        \bonne[A2]{OK}\scoring{b=1}
    }
    \explain{
    La pile zinc-air offre une autonomie de 91 h, contre 22,6 h pour la pile zinc-manganèse. Il est donc préférable de privilégier la pile zinc-air pour une plus grande autonomie.
    }
\end{questionmult}

\else
    %\section{Piles électrochimiques et alimentation d'une voiture télécommandée}
\textbf{Source du sujet~:} Baccalauréat STL Physique-Chimie, Polynésie – 2023.

\subsection{Introduction}
Une pile Daniell est un dispositif électrochimique réalisant la conversion d'énergie chimique
en énergie électrique par le biais d'une réaction d'oxydoréduction. Le principe de
fonctionnement de cette pile a été démontré par Edmond Becquerel et a été perfectionné
par le chimiste britannique John Daniell en 1836, au moment où le développement du
télégraphe faisait apparaître un besoin urgent de sources de tension sûres et constantes.

La pile Daniell vient corriger certains défauts de la pile Volta : elle est simple de
construction, commode d'usage et de tension stable durant son utilisation, si bien qu'elle
servit pendant longtemps de pile étalon en laboratoire.

On appelle pile alcaline un type de pile électrique dont l'électrolyte est alcalin, c’est-à-dire basique. L’un des modèles les plus courants est la pile alcaline zinc-dioxyde de manganèse (Zn-MnO2) que l’on utilise pour alimenter des jouets par exemple (\emph{Source
Wikipédia}).

L’objectif de cet exercice est d’étudier une pile Daniell réalisée au laboratoire puis d’étudier les performances d’une voiture télécommandée, alimentée à l’aide de piles alcalines.

\subsection{Données}
\begin{itemize}
    \item Masses molaires atomiques : $M(\mathrm{Cu}) = \SI{63.5}{\gram\per\mole}$ ; $M(\mathrm{Zn}) = \SI{65.4}{\gram\per\mole}$ ;
    \item Charge élémentaire : $e = 1,60 \times 10^{-19}\ \mathrm{C}$ ;
    \item Constante d'Avogadro : $N_{\mathrm{A}} = 6,02 \times 10^{23}\ \mathrm{mol}^{-1}$ ;
    \item Couples oxydant/réducteur de la pile Daniell : $\mathrm{Zn^{2+}(aq) / Zn(s)}$ et $\mathrm{Cu^{2+}(aq) / Cu(s)}$ ;
    \item Couples oxydant/réducteur de la pile alcaline : $\mathrm{ZnO(s) / Zn(s)}$ et $\mathrm{MnO_2(s) / MnO(OH)(s)}$ ;
    \item Capacité d'une pile alcaline : de 2 000 mAh à 3 000 mAh, soit respectivement de 7 200 C à 10 800 C ;
    \item Masse de la voiture radiocommandée : $m = \SI{741}{\gram}$.
\end{itemize}

\subsection{Étude d'une pile Daniell}
Une pile Daniell est constituée d'une lame de zinc plongée dans une solution contenant un volume $V_{0} = \SI{20}{\milli\liter}$ d'une solution de sulfate de zinc ($\mathrm{Zn^{2+}(aq) + SO_4^{2-}(aq)}$) de concentration $C_0 = \SI{1.0e-2}{\mole\per\liter}$ et d'une lame de cuivre plongée dans une solution contenant un volume $V_{0} = \SI{20}{\milli\liter}$ de solution de sulfate de cuivre ($\mathrm{Cu^{2+}(aq) + SO_4^{2-}(aq)}$) de concentration $C_0 = \SI{1.0e-2}{\mole\per\liter}$. Les deux solutions sont reliées par un pont salin contenant une solution de nitrate d'ammonium ($\mathrm{NH_4^+(aq) + NO_3^-(aq)}$).

Un voltmètre est branché aux bornes de la pile pour mesurer la valeur de la tension à vide $U_{0}$ lorsque la pile ne débite pas de courant électrique. Le voltmètre mesure une tension à vide positive~: $U_{0} = \SI{1.102}{\volt}$.

\begin{questionmult}{pileD1}
    Préciser, en justifiant la réponse, la polarité de chaque électrode.
    \begin{EnvQuadrillage}[NbCarreaux=10x4,Grille=Seyes,Marge=1]
    \end{EnvQuadrillage}
    \AMCOpen{}{
        \mauvaise[F]{NA}\scoring{b=0}
        \bonne[A1]{Polarité correcte}\scoring{b=1}
        \bonne[A2]{Justification correcte}\scoring{b=1}
    }
    \explain{
    La tension mesurée est positive, ce qui indique que l'électrode de cuivre est la borne positive et l'électrode de zinc la borne négative.
    }
\end{questionmult}

\begin{questionmult}{pileD2}
    Écrire les équations de demi-réaction électroniques se produisant à l'anode et à la cathode. En déduire la réaction d'oxydoréduction modélisant la transformation chimique au sein de la pile.
    \begin{EnvQuadrillage}[NbCarreaux=10x4,Grille=Seyes,Marge=1]
    \end{EnvQuadrillage}
    \AMCOpen{}{
        \mauvaise[F]{NA}\scoring{b=0}
        \bonne[A1]{Équations correctes}\scoring{b=1}
        \bonne[A2]{Réaction correcte}\scoring{b=1}
    }
    \explain{
    Les demi-équations sont :
    $$
    \ce{Zn(s) -> Zn^{2+}(aq) + 2 e^{-}}
    $$
    $$
    \ce{Cu^{2+}(aq) + 2 e^{-} -> Cu(s)}
    $$
    La réaction globale est :
    $$
    \ce{Zn(s) + Cu^{2+}(aq) -> Zn^{2+}(aq) + Cu(s)}
    $$
    }
\end{questionmult}

\begin{questionmult}{pileD3}
    Donner le sens de déplacement des ions ammonium $\mathrm{NH_4^+}$ (aq) et des ions nitrate $\mathrm{NO_3^-}$ (aq) dans le pont salin.
    \begin{EnvQuadrillage}[NbCarreaux=10x4,Grille=Seyes,Marge=1]
    \end{EnvQuadrillage}
    \AMCOpen{}{
        \mauvaise[F]{NA}\scoring{b=0}
        \bonne[A1]{Sens correct}\scoring{b=1}
    }
    \explain{
    Les ions ammonium $\mathrm{NH_4^+}$ se déplacent vers la solution de sulfate de cuivre, et les ions nitrate $\mathrm{NO_3^-}$ se déplacent vers la solution de sulfate de zinc.
    }
\end{questionmult}

\begin{questionmult}{pileD4}
    Déterminer le réactif limitant de la réaction.
    \begin{EnvQuadrillage}[NbCarreaux=10x4,Grille=Seyes,Marge=1]
    \end{EnvQuadrillage}
    \AMCOpen{}{
        \mauvaise[F]{NA}\scoring{b=0}
        \bonne[A1]{Réactif correct}\scoring{b=1}
        \bonne[A2]{Calcul correct}\scoring{b=1}
    }
    \explain{
    La quantité de matière initiale de zinc est $n_{\mathrm{Zn}} = \frac{m_{\mathrm{Zn}}}{M(\mathrm{Zn})} = \frac{50.2}{65.4} = \SI{0.767}{\mole}$.
    La quantité de matière initiale de cuivre est $n_{\mathrm{Cu}} = \frac{m_{\mathrm{Cu}}}{M(\mathrm{Cu})} = \frac{62.8}{63.5} = \SI{0.989}{\mole}$.
    Le réactif limitant est le zinc.
    }
\end{questionmult}

\begin{questionmult}{pileD5}
    En déduire la capacité $Q_D$ de la pile Daniell. Commenter la valeur obtenue.
    \begin{EnvQuadrillage}[NbCarreaux=10x4,Grille=Seyes,Marge=1]
    \end{EnvQuadrillage}
    \AMCOpen{}{
        \mauvaise[F]{NA}\scoring{b=0}
        \bonne[A1]{Calcul correct}\scoring{b=1}
        \bonne[A2]{Valeur correcte}\scoring{b=1}
    }
    \explain{
    La quantité d'électricité est donnée par $Q_D = n_{\mathrm{e^{-}}} \times F = 2 \times n_{\mathrm{Zn}} \times F = 2 \times 0.767 \times 96500 = \SI{1.48e5}{\coulomb}$.
    }
\end{questionmult}

\subsection{Étude d'une pile alcaline Power+}
Les piles alcalines sont fabriquées sous forme de cylindres et de boutons. Une pile cylindrique est contenue dans un tube d'acier, qui sert de collecteur du courant de la cathode.

\begin{questionmult}{pileA1}
    Identifier en justifiant la réponse, le réactif limitant de la pile alcaline.
    \begin{EnvQuadrillage}[NbCarreaux=10x4,Grille=Seyes,Marge=1]
    \end{EnvQuadrillage}
    \AMCOpen{}{
        \mauvaise[F]{NA}\scoring{b=0}
        \bonne[A1]{Réactif correct}\scoring{b=1}
        \bonne[A2]{Justification correcte}\scoring{b=1}
    }
    \explain{
    Le réactif limitant est le zinc, car il est entièrement consommé avant le dioxyde de manganèse.
    }
\end{questionmult}

\subsection{Alimentation d'une voiture radiocommandée et étude des performances}
Une voiture télécommandée est alimentée par un bloc de 6 piles alcalines AA Power+ étudiées précédemment.

\begin{questionmult}{voiture1}
    Déterminer la valeur de la durée d'utilisation maximale attendue de la voiture télécommandée. Commenter le résultat.
    \begin{EnvQuadrillage}[NbCarreaux=10x4,Grille=Seyes,Marge=1]
    \end{EnvQuadrillage}
    \AMCOpen{}{
        \mauvaise[F]{NA}\scoring{b=0}
        \bonne[A1]{Calcul correct}\scoring{b=1}
        \bonne[A2]{Valeur correcte}\scoring{b=1}
    }
    \explain{
    La durée d'utilisation maximale est donnée par $t = \frac{Q_B}{I} = \frac{10000}{0.6} = \SI{1.67e4}{\second} = \SI{4.64}{\hour}$.
    }
\end{questionmult}

\fi

\section{Comparaison des piles pour appareil auditif}
\textbf{Source du sujet~:} Baccalauréat STL Physique-Chimie, Centres étrangers – 2022.

\subsection{Introduction}
Pour alimenter un appareil auditif, deux types de piles boutons sont disponibles :
\begin{itemize}
    \item une pile bouton zinc-manganèse ;
    \item une pile bouton zinc-air.
\end{itemize}

Leurs caractéristiques sont les suivantes :

\begin{center}
\begin{tabular}{|c|c|c|c|c|}
\hline
Type de pile & Quantité d'électricité disponible & Masse de la pile & Tension & Intensité du courant \\
\hline
zinc-manganèse & à déterminer & 0,5 g & 1,5 V & 2,0 mA \\
\hline
zinc-air & 182 mA·h & 0,5 g & 1,5 V & 2,0 mA \\
\hline
\end{tabular}
\end{center}

L'objectif de l'exercice est de déterminer la pile à privilégier pour une plus grande autonomie de fonctionnement de l'appareil auditif.

\subsection{Données}
\begin{itemize}
    \item Masses molaires atomiques : $M(\mathrm{Zn}) = \SI{65}{\gram\per\mole}$ ; $M(\mathrm{Mn}) = \SI{55}{\gram\per\mole}$ ; $M(\mathrm{O}) = \SI{16}{\gram\per\mole}$ ; $M(\mathrm{H}) = \SI{1,0}{\gram\per\mole}$ ;
    \item Constante de Faraday : $F = \SI{9,65e4}{\coulomb\per\mole}$.
\end{itemize}

\subsection{Partie A~: Étude de la pile zinc-manganèse}
Dans une pile bouton zinc-manganèse, une électrode contient du zinc $\ce{Zn(s)}$ et l'autre du dioxyde de manganèse $\ce{MnO2(s)}$ et du carbone $\ce{C(s)}$. Le carbone n'intervient pas dans le couple de la demi-pile mais permet notamment la conduction du courant.\\
Afin de déterminer la polarité de la pile, on branche la borne COM d’un voltmètre à
l’électrode de carbone et la borne V à l’électrode de zinc, on mesure alors une tension
de \SI{-1.5}{\volt}.\\

\begin{questionmult}{pileA1}
    Déduire de la mesure de la tension à vide la polarité de chaque électrode en précisant le raisonnement.\\
    \begin{EnvQuadrillage}[NbCarreaux=18x3,Grille=Seyes,Marge=1]
    \end{EnvQuadrillage}
    \AMCOpen{}{
        \mauvaise[F]{NA}\scoring{b=0}
        \bonne[A1]{idOK}\scoring{b=1}
        \bonne[A2]{valOK}\scoring{b=1}
    }
    \explain{
    La tension mesurée est $-1,50\ \mathrm{V}$ lorsque la borne COM est branchée sur l'électrode de carbone et la borne V sur l'électrode de zinc. Cela indique que l'électrode de carbone est la borne positive et l'électrode de zinc la borne négative.
    }
\end{questionmult}

\begin{questionmult}{pileA2}
    Indiquer la polarité des bornes de la pile, le sens de déplacement des électrons ainsi que le sens du courant sur la figure.

    \includegraphics[width=0.6\linewidth]{ste7_pile_auditif_pile.png}

    \AMCOpen{}{
        \mauvaise[F]{NA}\scoring{b=0}
        \bonne[A1]{PolOK}\scoring{b=1}
        \bonne[A2]{SensOK}\scoring{b=1}
    }
    \explain{
    Les électrons se déplacent de l'électrode de zinc vers l'électrode de carbone, et le courant circule dans le sens opposé.
    }
\end{questionmult}

 \vspace{0.5cm}

Les demi-équations électrochimiques se produisant aux électrodes de la pile sont les suivantes :

$$
\ce{Zn(s) + 2 HO^{-}(aq) -> ZnO(s) + H2O(l) + 2 e^{-}}
$$

$$
\ce{MnO2(s) + H2O(l) + 1 e^{-} -> MnO2H(s) + HO^{-}(aq)}
$$

\begin{questionmult}{pileA3}
    Associer chaque demi-équation électrochimique à l'électrode correspondante et préciser si elle est le siège d'une oxydation ou d'une réduction.\\
    \begin{EnvQuadrillage}[NbCarreaux=18x4,Grille=Seyes,Marge=1]
    \end{EnvQuadrillage}
    \AMCOpen{}{
        \mauvaise{NA}\scoring{b=0}
        \bonne{assOK}\scoring{b=1}
        \bonne{redoxOK}\scoring{b=1}
    }
    \explain{
    La première demi-équation correspond à l'oxydation du zinc à l'anode. La deuxième demi-équation correspond à la réduction du dioxyde de manganèse à la cathode.
    }
\end{questionmult}

\begin{questionmult}{pileA4}
    Établir l'équation de la réaction modélisant le fonctionnement de la pile, en précisant la démarche.\\
    \begin{EnvQuadrillage}[NbCarreaux=18x8,Grille=Seyes,Marge=1]
    \end{EnvQuadrillage}
    \AMCOpen{}{
        \mauvaise{NA}\scoring{b=0}
        \bonne{idK}\scoring{b=1}
        \bonne{eqOK}\scoring{b=1}
    }
    \explain{
    L'équation de la réaction est obtenue en combinant les deux demi-équations, en tenant compte du nombre d'électrons échangés :
    $$
    \ce{2 Zn(s) + MnO2(s) + H2O(l) -> 2 ZnO(s) + MnO2H(s)}
    $$
    }
\end{questionmult}

\begin{questionmult}{pileA5}
    Sachant que la masse initiale de l'électrode en zinc est $m_{\mathrm{Zn}} = 55\ \mathrm{mg}$, déterminer la quantité de matière initiale de zinc $n_{\mathrm{Zn}}$ présente dans la pile.\\
    \begin{EnvQuadrillage}[NbCarreaux=18x3,Grille=Seyes,Marge=1]
    \end{EnvQuadrillage}
    \AMCOpen{}{
        \mauvaise[F]{NA}\scoring{b=0}
        \bonne[A1]{idOK}\scoring{b=1}
        \bonne[A2]{valOK}\scoring{b=1}
    }
    \explain{
    La quantité de matière initiale de zinc est donnée par $n_{\mathrm{Zn}} = \frac{m_{\mathrm{Zn}}}{M(\mathrm{Zn})} = \frac{55 \times 10^{-3}}{65} = 8,46 \times 10^{-4}\ \mathrm{mol}$.
    }
\end{questionmult}

\begin{questionmult}{pileA6}
    En admettant que le zinc est entièrement consommé et que le dioxyde de manganèse $\ce{MnO2(s)}$ est en excès, déduire la quantité de matière $n_{\mathrm{e}^{-}}$ d'électrons échangés au cours de la transformation chimique.\\
    \begin{EnvQuadrillage}[NbCarreaux=18x3,Grille=Seyes,Marge=1]
    \end{EnvQuadrillage}
    \AMCOpen{}{
        \mauvaise[F]{NA}\scoring{b=0}
        \bonne[A1]{idOK}\scoring{b=1}
        \bonne[A2]{valOK}\scoring{b=1}
    }
    \explain{
    Chaque mole de zinc libère 2 moles d'électrons. Donc, $n_{\mathrm{e}^{-}} = 2 \times n_{\mathrm{Zn}} = 1,69 \times 10^{-3}\ \mathrm{mol}$.
    }
\end{questionmult}

\begin{questionmult}{pileA7}
    Vérifier que la valeur de la quantité d'électricité disponible dans la pile est voisine de $Q_{1} = 164\ \mathrm{C}$.\\
    \begin{EnvQuadrillage}[NbCarreaux=10x2,Grille=Seyes,Marge=1]
    \end{EnvQuadrillage}
    \AMCOpen{}{
        \mauvaise[F]{NA}\scoring{b=0}
        \bonne[A1]{idOK}\scoring{b=1}
        \bonne[A2]{valOK}\scoring{b=1}
    }
    \explain{
    La quantité d'électricité est donnée par $Q = n_{\mathrm{e}^{-}} \times F = 1,69 \times 10^{-3} \times 9,65 \times 10^{4} = 163\ \mathrm{C}$.
    }
\end{questionmult}

\begin{questionmult}{pileA8}
    Calculer la durée de fonctionnement de la pile zinc-manganèse.\\
    \begin{EnvQuadrillage}[NbCarreaux=10x2,Grille=Seyes,Marge=1]
    \end{EnvQuadrillage}
    \AMCOpen{}{
        \mauvaise[F]{NA}\scoring{b=0}
        \bonne[A1]{idOK}\scoring{b=1}
        \bonne[A2]{valOK}\scoring{b=1}
    }
    \explain{
    La durée de fonctionnement est donnée par $t = \frac{Q}{I} = \frac{163}{2,0 \times 10^{-3}} = 81\,500\ \mathrm{s} = 22,6\ \mathrm{h}$.
    }
\end{questionmult}

\subsection{Partie B~: Étude de la pile zinc-air}
L'équation de la réaction modélisant le fonctionnement de la pile correspond à l'oxydation du zinc en oxyde de zinc par le dioxygène de l'air :

$$
\ce{2 Zn(s) + O2(g) -> 2 ZnO(s)}
$$

Un des deux couples redox mis en jeu est le couple $\ce{O2(g) / H2O(l)}$.

\begin{questionmult}{pileB1}
    Écrire l'équation de demi-réaction électrochimique correspondant à ce couple et justifier le caractère oxydant du dioxygène.\\
    \begin{EnvQuadrillage}[NbCarreaux=18x4,Grille=Seyes,Marge=1]
    \end{EnvQuadrillage}
    \AMCOpen{}{
        \mauvaise[F]{NA}\scoring{b=0}
        \bonne[A1]{Équation correcte}\scoring{b=1}
        \bonne[A2]{Justification correcte}\scoring{b=1}
    }
    \explain{
    La demi-équation est $\ce{O2(g) + 2 H2O(l) + 4 e^{-} -> 4 HO^{-}(aq)}$. Le dioxygène est un oxydant car il gagne des électrons.
    }
\end{questionmult}

\begin{questionmult}{pileB2}
    Calculer la durée de fonctionnement de la pile zinc-air en vous aidant de ses caractéristiques.\\
    \begin{EnvQuadrillage}[NbCarreaux=18x4,Grille=Seyes,Marge=1]
    \end{EnvQuadrillage}
    \AMCOpen{}{
        \mauvaise[F]{NA}\scoring{b=0}
        \bonne[A1]{idOK}\scoring{b=1}
        \bonne[A2]{valOK}\scoring{b=1}
    }
    \explain{
    La quantité d'électricité disponible est $Q = 182\ \mathrm{mA\cdot h} = 655,2\ \mathrm{C}$. La durée de fonctionnement est $t = \frac{Q}{I} = \frac{655,2}{2,0 \times 10^{-3}} = 327\,600\ \mathrm{s} = 91\ \mathrm{h}$.
    }
\end{questionmult}

\subsection{Partie C~: Choix de la pile à privilégier}
\begin{questionmult}{pileC1}
    À l'aide des parties A et B et des informations disponibles, indiquer le type de pile qu'il faudrait privilégier pour l'alimentation d'un appareil auditif. Justifier la réponse.\\
    \begin{EnvQuadrillage}[NbCarreaux=18x4,Grille=Seyes,Marge=1]
    \end{EnvQuadrillage}
    \AMCOpen{}{
        \mauvaise[F]{NA}\scoring{b=0}
        \bonne[A1]{id}\scoring{b=1}
        \bonne[A2]{OK}\scoring{b=1}
    }
    \explain{
    La pile zinc-air offre une autonomie de 91 h, contre 22,6 h pour la pile zinc-manganèse. Il est donc préférable de privilégier la pile zinc-air pour une plus grande autonomie.
    }
\end{questionmult}

\section{Piles électrochimiques et alimentation d'une voiture télécommandée}
\textbf{Source du sujet~:} Baccalauréat STL Physique-Chimie, Polynésie – 2023.

\subsection{Introduction}
Une pile Daniell est un dispositif électrochimique réalisant la conversion d'énergie chimique
en énergie électrique par le biais d'une réaction d'oxydoréduction. Le principe de
fonctionnement de cette pile a été démontré par Edmond Becquerel et a été perfectionné
par le chimiste britannique John Daniell en 1836, au moment où le développement du
télégraphe faisait apparaître un besoin urgent de sources de tension sûres et constantes.

La pile Daniell vient corriger certains défauts de la pile Volta : elle est simple de
construction, commode d'usage et de tension stable durant son utilisation, si bien qu'elle
servit pendant longtemps de pile étalon en laboratoire.

On appelle pile alcaline un type de pile électrique dont l'électrolyte est alcalin, c’est-à-dire basique. L’un des modèles les plus courants est la pile alcaline zinc-dioxyde de manganèse (Zn-MnO2) que l’on utilise pour alimenter des jouets par exemple (\emph{Source
Wikipédia}).

L’objectif de cet exercice est d’étudier une pile Daniell réalisée au laboratoire puis d’étudier les performances d’une voiture télécommandée, alimentée à l’aide de piles alcalines.

\subsection{Données}
\begin{itemize}
    \item Masses molaires atomiques : $M(\mathrm{Cu}) = \SI{63.5}{\gram\per\mole}$ ; $M(\mathrm{Zn}) = \SI{65.4}{\gram\per\mole}$ ;
    \item Charge élémentaire : $e = 1,60 \times 10^{-19}\ \mathrm{C}$ ;
    \item Constante d'Avogadro : $N_{\mathrm{A}} = 6,02 \times 10^{23}\ \mathrm{mol}^{-1}$ ;
    \item Couples oxydant/réducteur de la pile Daniell : $\mathrm{Zn^{2+}(aq) / Zn(s)}$ et $\mathrm{Cu^{2+}(aq) / Cu(s)}$ ;
    \item Couples oxydant/réducteur de la pile alcaline : $\mathrm{ZnO(s) / Zn(s)}$ et $\mathrm{MnO_2(s) / MnO(OH)(s)}$ ;
    \item Capacité d'une pile alcaline : de 2 000 mAh à 3 000 mAh, soit respectivement de 7 200 C à 10 800 C ;
    \item Masse de la voiture radiocommandée : $m = \SI{741}{\gram}$.
\end{itemize}

\subsection{Étude d'une pile Daniell}
Une pile Daniell est constituée d'une lame de zinc plongée dans une solution contenant un volume $V_{0} = \SI{20}{\milli\liter}$ d'une solution de sulfate de zinc ($\mathrm{Zn^{2+}(aq) + SO_4^{2-}(aq)}$) de concentration $C_0 = \SI{1.0e-2}{\mole\per\liter}$ et d'une lame de cuivre plongée dans une solution contenant un volume $V_{0} = \SI{20}{\milli\liter}$ de solution de sulfate de cuivre ($\mathrm{Cu^{2+}(aq) + SO_4^{2-}(aq)}$) de concentration $C_0 = \SI{1.0e-2}{\mole\per\liter}$. Les deux solutions sont reliées par un pont salin contenant une solution de nitrate d'ammonium ($\mathrm{NH_4^+(aq) + NO_3^-(aq)}$).

Un voltmètre est branché aux bornes de la pile pour mesurer la valeur de la tension à vide $U_{0}$ lorsque la pile ne débite pas de courant électrique. Le voltmètre mesure une tension à vide positive~: $U_{0} = \SI{1.102}{\volt}$.

\begin{questionmult}{pileD1}
    Préciser, en justifiant la réponse, la polarité de chaque électrode.
    \begin{EnvQuadrillage}[NbCarreaux=10x4,Grille=Seyes,Marge=1]
    \end{EnvQuadrillage}
    \AMCOpen{}{
        \mauvaise[F]{NA}\scoring{b=0}
        \bonne[A1]{Polarité correcte}\scoring{b=1}
        \bonne[A2]{Justification correcte}\scoring{b=1}
    }
    \explain{
    La tension mesurée est positive, ce qui indique que l'électrode de cuivre est la borne positive et l'électrode de zinc la borne négative.
    }
\end{questionmult}

\begin{questionmult}{pileD2}
    Écrire les équations de demi-réaction électroniques se produisant à l'anode et à la cathode. En déduire la réaction d'oxydoréduction modélisant la transformation chimique au sein de la pile.
    \begin{EnvQuadrillage}[NbCarreaux=10x4,Grille=Seyes,Marge=1]
    \end{EnvQuadrillage}
    \AMCOpen{}{
        \mauvaise[F]{NA}\scoring{b=0}
        \bonne[A1]{Équations correctes}\scoring{b=1}
        \bonne[A2]{Réaction correcte}\scoring{b=1}
    }
    \explain{
    Les demi-équations sont :
    $$
    \ce{Zn(s) -> Zn^{2+}(aq) + 2 e^{-}}
    $$
    $$
    \ce{Cu^{2+}(aq) + 2 e^{-} -> Cu(s)}
    $$
    La réaction globale est :
    $$
    \ce{Zn(s) + Cu^{2+}(aq) -> Zn^{2+}(aq) + Cu(s)}
    $$
    }
\end{questionmult}

\begin{questionmult}{pileD3}
    Donner le sens de déplacement des ions ammonium $\mathrm{NH_4^+}$ (aq) et des ions nitrate $\mathrm{NO_3^-}$ (aq) dans le pont salin.
    \begin{EnvQuadrillage}[NbCarreaux=10x4,Grille=Seyes,Marge=1]
    \end{EnvQuadrillage}
    \AMCOpen{}{
        \mauvaise[F]{NA}\scoring{b=0}
        \bonne[A1]{Sens correct}\scoring{b=1}
    }
    \explain{
    Les ions ammonium $\mathrm{NH_4^+}$ se déplacent vers la solution de sulfate de cuivre, et les ions nitrate $\mathrm{NO_3^-}$ se déplacent vers la solution de sulfate de zinc.
    }
\end{questionmult}

\begin{questionmult}{pileD4}
    Déterminer le réactif limitant de la réaction.
    \begin{EnvQuadrillage}[NbCarreaux=10x4,Grille=Seyes,Marge=1]
    \end{EnvQuadrillage}
    \AMCOpen{}{
        \mauvaise[F]{NA}\scoring{b=0}
        \bonne[A1]{Réactif correct}\scoring{b=1}
        \bonne[A2]{Calcul correct}\scoring{b=1}
    }
    \explain{
    La quantité de matière initiale de zinc est $n_{\mathrm{Zn}} = \frac{m_{\mathrm{Zn}}}{M(\mathrm{Zn})} = \frac{50.2}{65.4} = \SI{0.767}{\mole}$.
    La quantité de matière initiale de cuivre est $n_{\mathrm{Cu}} = \frac{m_{\mathrm{Cu}}}{M(\mathrm{Cu})} = \frac{62.8}{63.5} = \SI{0.989}{\mole}$.
    Le réactif limitant est le zinc.
    }
\end{questionmult}

\begin{questionmult}{pileD5}
    En déduire la capacité $Q_D$ de la pile Daniell. Commenter la valeur obtenue.
    \begin{EnvQuadrillage}[NbCarreaux=10x4,Grille=Seyes,Marge=1]
    \end{EnvQuadrillage}
    \AMCOpen{}{
        \mauvaise[F]{NA}\scoring{b=0}
        \bonne[A1]{Calcul correct}\scoring{b=1}
        \bonne[A2]{Valeur correcte}\scoring{b=1}
    }
    \explain{
    La quantité d'électricité est donnée par $Q_D = n_{\mathrm{e^{-}}} \times F = 2 \times n_{\mathrm{Zn}} \times F = 2 \times 0.767 \times 96500 = \SI{1.48e5}{\coulomb}$.
    }
\end{questionmult}

\subsection{Étude d'une pile alcaline Power+}
Les piles alcalines sont fabriquées sous forme de cylindres et de boutons. Une pile cylindrique est contenue dans un tube d'acier, qui sert de collecteur du courant de la cathode.

\begin{questionmult}{pileA1}
    Identifier en justifiant la réponse, le réactif limitant de la pile alcaline.
    \begin{EnvQuadrillage}[NbCarreaux=10x4,Grille=Seyes,Marge=1]
    \end{EnvQuadrillage}
    \AMCOpen{}{
        \mauvaise[F]{NA}\scoring{b=0}
        \bonne[A1]{Réactif correct}\scoring{b=1}
        \bonne[A2]{Justification correcte}\scoring{b=1}
    }
    \explain{
    Le réactif limitant est le zinc, car il est entièrement consommé avant le dioxyde de manganèse.
    }
\end{questionmult}

\subsection{Alimentation d'une voiture radiocommandée et étude des performances}
Une voiture télécommandée est alimentée par un bloc de 6 piles alcalines AA Power+ étudiées précédemment.

\begin{questionmult}{voiture1}
    Déterminer la valeur de la durée d'utilisation maximale attendue de la voiture télécommandée. Commenter le résultat.
    \begin{EnvQuadrillage}[NbCarreaux=10x4,Grille=Seyes,Marge=1]
    \end{EnvQuadrillage}
    \AMCOpen{}{
        \mauvaise[F]{NA}\scoring{b=0}
        \bonne[A1]{Calcul correct}\scoring{b=1}
        \bonne[A2]{Valeur correcte}\scoring{b=1}
    }
    \explain{
    La durée d'utilisation maximale est donnée par $t = \frac{Q_B}{I} = \frac{10000}{0.6} = \SI{1.67e4}{\second} = \SI{4.64}{\hour}$.
    }
\end{questionmult}


\AMCcleardoublepage
}

\csvreader[head to column names]{liste.csv}{}{\sujet}

\end{document}
